\begin{workedexample}[Worked Example: Fields and power in a plane wave]
A uniform plane wave in free space has a peak electric field
$E_0 = 10\ \mathrm{V/m}$. The magnetic field follows from the intrinsic impedance,
\[ H_0 = \frac{E_0}{\eta_0} = \frac{10}{376.7} = 2.65\times10^{-2}\ \mathrm{A/m}. \]
The time-average power density carried by the wave is
\[ S_{\text{avg}} = \frac{E_0^{2}}{2\eta_0} = \frac{10^{2}}{2(376.7)}
   = 0.133\ \mathrm{W/m^{2}}. \]
$\mathbf{E}$, $\mathbf{H}$, and the direction of propagation are mutually
perpendicular, and the power flows along $\mathbf{S} = \mathbf{E}\times\mathbf{H}$.
\end{workedexample}

\begin{keytakeaways}
\begin{itemize}[leftmargin=1.4em,itemsep=2pt]
\item In a TEM wave $\mathbf{E}\perp\mathbf{H}\perp$ direction of travel, and
$|E|/|H| = \eta$ (the intrinsic impedance).
\item Free space: $\eta_0 = 120\pi \approx 377\ \Omega$; in a medium
$\eta = \eta_0\sqrt{\mu_r/\varepsilon_r}$.
\item Time-average power density $S_{\text{avg}} = E_0^{2}/(2\eta)$; polarization is
the tip-path traced by $\mathbf{E}$ over one cycle.
\end{itemize}
\end{keytakeaways}

\begin{exercises}
\item A wave in free space has a power density of $1\ \mathrm{mW/cm^{2}}$
($=10\ \mathrm{W/m^{2}}$). Find its peak electric field.
  \begin{answer}\textbf{Answer.} $E_0 = \sqrt{2\eta_0 S} = \sqrt{2(376.7)(10)} = \sqrt{7534} = 86.8\ \mathrm{V/m}$.\end{answer}
\item What is the free-space wavelength at $2.45\ \mathrm{GHz}$?
  \begin{answer}\textbf{Answer.} $\lambda = c/f = (3\times10^{8})/(2.45\times10^{9}) = 0.1224\ \mathrm{m} = 12.2\ \mathrm{cm}$.\end{answer}
\item Two equal-amplitude, orthogonal linear components combine to give circular
polarization. What phase difference between them is required?
  \begin{answer}\textbf{Answer.} $90^{\circ}$ (a quarter cycle); the handedness depends on the sign.\end{answer}
\end{exercises}
